\scanpage{27}
\mspage{7}
à \inserted{différentielles des} op. différentiels.

Dans b) et c) reste posée, essentiellement, la question de la définition de la bonne \uline{catégorie de coefficients}, pour englober à la fois le formalisme des complexes \deleted{quasi} ``linéaires'' cohérents, et les complexes du type De Rham.
Les exposés que je fais doivent donner quelques indications plus précises sur la voie \deleted{\unsure{devons}}\deleted{\unsure{suivre}} dans laquelle il conviendra \deleted{de} probablement de chercher ces catégories de coefficients.

§ 2 \uline{Tour d'horizon sur les \deleted{autres} théories cohomologiques \inserted{$\ell$-adiques} pour les variétés algébriques}

\uline{a) Avantages}. \deleted{\unsure{Inconvénients}.}

On peut évidemment se demander : pourquoi à tout prix \deleted{désirer} aller jusqu'au bout dans un programme d'algébrisation \inserted{et de généralisation} à outrance de la cohomologie de De Rham.
On dispose déjà des invariants cohomologiques $\ell$-adiques, qui se prêtent à des applications \deleted{géométriques}
\margin{26}
